% dilworth.tex

%%%%%%%%%%%%%%%
\begin{frame}{}
  Dilworth
\end{frame}
%%%%%%%%%%%%%%%

%%%%%%%%%%%%%%%
\begin{frame}{}
  Equivalent theorems
\end{frame}
%%%%%%%%%%%%%%%

%%%%%%%%%%%%%%%
\begin{frame}{}
  \begin{definition}[可比与不可比 (Comparable/Incomparable)]
    设 $(X, \preceq)$ 是偏序集。
    对 $a, b \in X$,

    \begin{description}
      \item[可比的 (Comparable):]
        \[
          a \preceq b \lor b \preceq a
        \]
      \item[不可比的 (Incomparable):]
        \[
          \lnot(a \preceq b \lor b \preceq a)
        \]
    \end{description}
  \end{definition}

  \fig{width = 0.50\textwidth}{figs/Hasse-Powerset}
\end{frame}
%%%%%%%%%%%%%%%

%%%%%%%%%%%%%%%
\begin{frame}{}
  \begin{definition}[链与反链 (Chain; Antichain)]
    设 $(X, \preceq)$ 是偏序集。

    \begin{itemize}
      \setlength{\itemsep}{6pt}
      \item 设 $S \subseteq X$ 且 $S$ 中元素两两\blue{可比}, 则称 $S$ 是\red{链}。
      \item 设 $S \subseteq X$ 且 $S$ 中元素两两\blue{不可比}, 则称 $S$ 是\red{反链}。
    \end{itemize}
  \end{definition}

  \begin{center}
    注意: 单元素集合既是链, 也是反链
  \end{center}

  \fig{width = 0.50\textwidth}{figs/Hasse-Powerset}
\end{frame}
%%%%%%%%%%%%%%%

%%%%%%%%%%%%%%%
\begin{frame}{}
  \begin{definition}[链分解 (Chain Decomposition)]
    设 $(X, \preceq)$ 是偏序集。
    对于 $C_{1}, C_{2}, \dots, C_{k} \subseteq X$, 如果
    \begin{enumerate}
      \setlength{\itemsep}{6pt}
      \item 每个 $C_{i}$ 都是\red{\bf 链};
      \item $C_{1}, C_{2}, \dots, C_{k}$ \red{\bf 互不相交};
      \item $X = \bigcup_{i=1}^{k} C_{i}$,
    \end{enumerate}
    则称 $\set{C_{1}, C_{2}, \dots, C_{k}}$
    是 $X$ 的\blue{\bf 链分解}。
  \end{definition}

  \fig{width = 0.40\textwidth}{figs/Hasse-Powerset}
\end{frame}
%%%%%%%%%%%%%%%

%%%%%%%%%%%%%%%
\begin{frame}{}
  \[
    \uncover<2->{\set{\blue{\set{x}}, \red{\set{y}}, \cyan{\set{z}}}}
    \qquad
    \uncover<3->{\set{\teal{\set{x, y}}, \brown{\set{y, z}}, \violet{\set{x, z}}}}
  \]
  \fig{width = 0.50\textwidth}{figs/Hasse-Powerset}
  \[
    \uncover<4->{\set{\blue{\set{\emptyset, \set{x}, \set{x, y}, \set{x, y, z}}},
      \red{\set{\set{y}, \set{y, z}}}, \cyan{\set{\set{z}, \set{x, z}}}}}
  \]
\end{frame}
%%%%%%%%%%%%%%%

%%%%%%%%%%%%%%%
\begin{frame}{}
  \fig{width = 0.50\textwidth}{figs/lattice-integer}
\end{frame}
%%%%%%%%%%%%%%%

%%%%%%%%%%%%%%%
\begin{frame}{}
  \fig{width = 0.70\textwidth}{figs/division-6}
\end{frame}
%%%%%%%%%%%%%%%

%%%%%%%%%%%%%%%
\begin{frame}{}
  \fig{width = 0.60\textwidth}{figs/powerset}
\end{frame}
%%%%%%%%%%%%%%%

%%%%%%%%%%%%%%%
\begin{frame}{}
  \fig{width = 0.70\textwidth}{figs/60-poset}
\end{frame}
%%%%%%%%%%%%%%%

%%%%%%%%%%%%%%%
\begin{frame}{}
  \fig{width = 0.60\textwidth}{figs/set-partition-refinement}
\end{frame}
%%%%%%%%%%%%%%%

%%%%%%%%%%%%%%%
\begin{frame}{}
  \begin{theorem}[Dilworth's Theorem]
    \red{最大}\blue{反链}的大小 = \red{最小}\blue{链分解}中链的条数
  \end{theorem}
\end{frame}
%%%%%%%%%%%%%%%

%%%%%%%%%%%%%%%
\begin{frame}{}
  这将是迄今为止我们在本课程中证明的最难的定理。

  \begin{columns}
    \column{0.50\textwidth}
    \column{0.50\textwidth}
      精密仪器、环环相扣
  \end{columns}
\end{frame}
%%%%%%%%%%%%%%%

%%%%%%%%%%%%%%%
\begin{frame}{}
  \begin{theorem}[Dilworth's Theorem]
    \red{最小}\blue{链分解}中链的条数 $=$ \red{最大}\blue{反链}的大小
  \end{theorem}

  \uncover<3->{
    \fig{width = 0.30\textwidth}{figs/so-easy}
  }

  \pause
  \begin{lemma}[$\le$: 链分解条数的\green{下界}]
    \red{最小}\blue{链分解}中链的条数 $\ge$ \red{最大}\blue{反链}的大小
  \end{lemma}
\end{frame}
%%%%%%%%%%%%%%%

%%%%%%%%%%%%%%%
\begin{frame}{}
  \begin{lemma}[$\le$: 链分解条数的\green{上界}]
    \red{最小}\blue{链分解}中链的条数 $\le$ \red{最大}\blue{反链}的大小
  \end{lemma}

  \pause
  \begin{center}
    ``这么多'' 就够了

    \pause
    假设 \red{最大}\blue{反链}的大小为 $k$

    \pause
    只需证明存在一个\red{链分解}包含 $k$ 条链

    \pause
    \blue{\large \bf 构造性证明:}
    由给定的大小为 $k$ 的反链构造 $k$ 条链

    \pause
    每条链包含反链中的一个元素

    \pause
    链的拼接 + 图示
  \end{center}
\end{frame}
%%%%%%%%%%%%%%%

%%%%%%%%%%%%%%%
\begin{frame}{}
  \begin{lemma}[$\le$: 链分解条数的\green{上界}]
    \red{最小}\blue{链分解}中链的条数 $\le$ \red{最大}\blue{反链}的大小
  \end{lemma}

  \begin{proof}
    \pause
    \blue{\bf 数学归纳法: \pause 对偏序集的大小 $|X|$ 进行归纳}

    \pause
    \vspace{1em}
    \begin{description}
      \setlength{\itemsep}{8pt}
      \item[基础情况:] $|X| = 1$, 定理显然成立
      \pause
      \item[归纳假设:] 假设对大小 $|X| < n$ 的偏序集, 定理均成立
      \pause
      \item[归纳步骤:] 考察大小 $|X| = n$ 的偏序集 $(X, \preceq)$
        \vspace{0.20cm}
        \begin{itemize}
          \setlength{\itemsep}{6pt}
          \pause
          \item 设 $X$ 的\red{最大}\blue{反链}的大小为 $k$
          \pause
          \item 需要构造一个包含 $k$ 条链的\red{链分解}
        \end{itemize}
    \end{description}
  \end{proof}
\end{frame}
%%%%%%%%%%%%%%%

%%%%%%%%%%%%%%%
\begin{frame}{}
\end{frame}
%%%%%%%%%%%%%%%

%%%%%%%%%%%%%%%
\begin{frame}{}
\end{frame}
%%%%%%%%%%%%%%%

%%%%%%%%%%%%%%%
\begin{frame}{}
\end{frame}
%%%%%%%%%%%%%%%

%%%%%%%%%%%%%%%
\begin{frame}{}
\end{frame}
%%%%%%%%%%%%%%%